% Unit 100 English translation of chapter7.tex lines 2141--2305.
% Source: Methods of Algebra, Volume 2, commit
% 9a5803ff77dd3257484cb177f851a73770a59dd3 (CC BY 4.0).
% stable-unit-id: o014.aljabr2.chapter7.module-descent
% Coverage: complete Section 7.9, ``Application: Descent of Modules'';
% stops immediately before Section 7.10 at line 2306.

% segment-id: o014.li2.chapter7.module-descent.h001
\section{Application: Descent of Modules}\label{sec:descent-modules}
\hypertarget{o014.aljabr2.chapter7.module-descent}{ }
% segment-id: o014.li2.chapter7.module-descent.p001
In geometry, given an adjoint pair
% segment-id: o014.li2.chapter7.module-descent.q001
$\begin{tikzcd}
	F: \mathcal{C} \arrow[shift left, r] & \mathcal{D}: G \arrow[shift left, l]
\end{tikzcd}$,
the process of reconstructing information on $\mathcal{C}$ from comodules on
$\mathcal{D}$ under the coaction of the corresponding comonad is usually
called \emph{descent}. This section explains the case of module theory.
\index{descent}

% segment-id: o014.li2.chapter7.module-descent.p002
Consider a homomorphism of commutative rings $K\to L$. The corresponding
forgetful functor $\mathcal{F}_{L|K}:L\dcate{Mod}\to K\dcate{Mod}$ fits into
the adjoint pair
% segment-id: o014.li2.chapter7.module-descent.q002
\[\begin{tikzcd}
	L \dotimes{K} (\cdot): K\dcate{Mod} \arrow[shift left, r] & L\dcate{Mod} : \mathcal{F}_{L|K}, \arrow[shift left, l]
\end{tikzcd}\]
We apply the results from the second half of \S\ref{sec:Morita} to study the
corresponding comonad on $L\dcate{Mod}$. More precisely, take $\Bbbk=K$ and
apply the ``change of rings'' construction of
Example~\ref{eg:comonad-change-ring}. Apart from notation, we use left rather
than right modules here. Since only commutative rings are considered, this
makes no substantive difference.

% segment-id: o014.li2.chapter7.module-descent.p003
Following the description in Example~\ref{eg:comonad-change-ring}, denote the
comonad determined on $L\dcate{Mod}$ by the adjoint pair by
$(\mathbf{L},\delta,\epsilon)$, where
% segment-id: o014.li2.chapter7.module-descent.q003
\[ \mathbf{L} = L \dotimes{K} \mathcal{F}_{L|K}(\cdot), \quad \delta: \mathbf{L} \to \mathbf{L}^2, \quad \epsilon: \mathbf{L} \to \identity. \]
From now on, as usual, the forgetful functor $\mathcal{F}_{L|K}$ will be
omitted from formulas to simplify the notation.

% segment-id: o014.li2.chapter7.module-descent.p004
A comodule over the comonad $\mathbf{L}$ is a datum $(M,a)$, where $M$ is an
$L$-module and $a:M\to L\dotimes{K}M$ is an $L$-module homomorphism, subject
to the requirement that the following diagrams commute:
% segment-id: o014.li2.chapter7.module-descent.q004
\[\begin{tikzcd}
	\ell \otimes 1 \otimes m \arrow[phantom, r, "\in" description] & L \dotimes{K} L \dotimes{K} M & L \dotimes{K} M \arrow[l, "{\mathbf{L}a = \identity_L \otimes a}"' inner sep=0.6em] \\
	\ell \otimes m \arrow[phantom, r, "\in" description] \arrow[mapsto, u] & L \dotimes{K} M \arrow[u, "\delta_M"] & M \arrow[l, "a"] \arrow[u, "a"']
\end{tikzcd} \quad \begin{tikzcd}
	\ell m \arrow[phantom, d, "\in" description, sloped] & \ell \otimes m \arrow[phantom, d, "\in" description, sloped] \arrow[mapsto, l] \\
	M & L \dotimes{K} M \arrow[l, "\epsilon_M"'] \\
	& M. \arrow[lu, "\identity_M"] \arrow[u, "a"']
\end{tikzcd}\]

% segment-id: o014.li2.chapter7.module-descent.p005
All comodules $(M,a)$ over the comonad $\mathbf{L}$ form a category
$L\dcate{Mod}^{\mathbf{L}}$. The dual of Lemma~\ref{prop:T-comparison}
gives the factorization
% segment-id: o014.li2.chapter7.module-descent.q005
\[ L \dotimes{K} (\cdot) = \left[ K\dcate{Mod} \xrightarrow{\mathbb{K}} L\dcate{Mod}^{\mathbf{L}} \xrightarrow{U^{\mathbf{L}}} L\dcate{Mod} \right] \; \text{as a composite}. \]
The functor $U^{\mathbf{L}}$ forgets the coaction, $(M,a)\mapsto M$, while
$\mathbb{K}$ maps an object $N$ to $(L\dotimes{K}N,a)$, where
% segment-id: o014.li2.chapter7.module-descent.q006
\begin{equation}\label{eqn:N-comonad-action}
	\begin{aligned}
		a: L \dotimes{K} N & \to L \dotimes{K} (L \dotimes{K} N) \quad \text{($L$-module homomorphism)} \\
		\ell \otimes x & \mapsto \ell \otimes 1 \otimes x.
	\end{aligned}
\end{equation}

% segment-id: o014.li2.chapter7.module-descent.p006
The central question is whether the adjoint pair is comonadic; equivalently,
is $\mathbb{K}$ an equivalence?

% segment-id: o014.li2.chapter7.module-descent.d001
\begin{definition}
	\index{faithfully flat}
	Given a homomorphism of commutative rings $K\to L$, if
	$L\dotimes{K}(\mathord\cdot):K\dcate{Mod}\to L\dcate{Mod}$ is a
	faithful exact functor, then $L$, as a commutative $K$-algebra, is called
	\emph{faithfully flat}.
\end{definition}

% segment-id: o014.li2.chapter7.module-descent.th001
\begin{theorem}[Flat descent]\label{prop:flat-descent}
	\index{descent!flat}
	Let $L$ be a faithfully flat commutative $K$-algebra. Then the adjoint pair
	$\left(L\dotimes{K}(\mathord\cdot),\mathcal{F}_{L|K}\right)$ is
	comonadic. In this case, a quasi-inverse $\mathbb{L}$ to
	$K\dcate{Mod}\xrightarrow{\mathbb{K}}L\dcate{Mod}^{\mathbf{L}}$ is
	described by the following equalizer diagram in $K\dcate{Mod}$:
	% segment-id: o014.li2.chapter7.module-descent.q007
	\[\begin{tikzcd}[row sep=tiny]
		\mathbb{L}(M, a) \arrow[r] & M \arrow[shift left, r, "a"] \arrow[shift right, r, "{m \mapsto 1 \otimes m}"'] & L \dotimes{K} M ,
	\end{tikzcd} \quad (M, a) \in \Obj\left(L\dcate{Mod}^{\mathbf{L}}\right). \]
\end{theorem}
% segment-id: o014.li2.chapter7.module-descent.pf001
\begin{proof}
	Proposition~\ref{prop:Morita-descent} establishes comonadicity, while the
	dual of Theorem~\ref{prop:Beck} already contains the description of the
	quasi-inverse functor.
\end{proof}

% segment-id: o014.li2.chapter7.module-descent.r001
\begin{remark}
	It must be emphasized that at the level of the derived categories
	$\cate{D}(K\dcate{Mod})$ and $\cate{D}(L\dcate{Mod})$, although the
	adjoint pair in question also has a derived version as in
	Example~\ref{eg:change-of-ring-adjunction} and
	Example~\ref{eg:change-of-rings}, there is no corresponding descent
	theorem.
\end{remark}

% segment-id: o014.li2.chapter7.module-descent.p007
The proof of flat descent is surprisingly simple. In practice, however, the
comonad or the descent data must be described in a usable form; this is the
theme of \S\ref{sec:descent-Galois}. In the remainder of this section we
recast the definition of $L\dcate{Mod}^{\mathbf{L}}$ in another form for
easy comparison with the literature. In essence, this is a concrete
rewriting of Proposition~\ref{prop:Morita-comonad} (in its left-module
version, with $P=L=P^\vee$), though the procedure is somewhat circuitous.

% segment-id: o014.li2.chapter7.module-descent.p008
For $i=1,2$, write $\iota_i:L\to L\dotimes{K}L$ for the homomorphism that
inserts an element into the $i$th tensor slot. For every $L$-module $M$ there
is a canonical isomorphism of $L$-modules
% segment-id: o014.li2.chapter7.module-descent.q008
\begin{align*}
	L \dotimes{K} M & \rightiso (L \dotimes{K} L) \dotimes{L, \iota_2} M \\
	\ell \otimes m & \mapsto \ell \otimes 1 \otimes m,
\end{align*}
where the right-hand side is made an $L$-module through $\iota_1$. By the
standard adjunction in module theory, specifying an $L$-module homomorphism
$a:M\to L\dotimes{K}M$ is therefore equivalent to specifying an
$L\dotimes{K}L$-module homomorphism
% segment-id: o014.li2.chapter7.module-descent.q009
\begin{equation}\label{eqn:flat-descent-map-aux}\begin{aligned}
		\tilde{a}: (L \dotimes{K} L) \dotimes{L, \iota_1} M & \to (L \dotimes{K} L) \dotimes{L, \iota_2} M \\
		1 \otimes 1 \otimes m & \mapsto \sum_{i=1}^n \ell_i \otimes 1 \otimes m_i ;
\end{aligned}\end{equation}
here $a(m)=\sum_{i=1}^n\ell_i\otimes m_i$. Since
$\iota_1|_K=\iota_2|_K$, when $M$ comes from a $K$-module $N$, careful
substitution from \eqref{eqn:N-comonad-action} shows that $\tilde a$ becomes
$\identity:L\dotimes{K}L\dotimes{K}N\to
L\dotimes{K}L\dotimes{K}N$.

% segment-id: o014.li2.chapter7.module-descent.p009
For brevity, write $L^{\otimes n}:=L\dotimes{K}\cdots\dotimes{K}L$ (with
$n$ factors). Temporarily forget $L\dcate{Mod}^{\mathbf{L}}$ and consider
an arbitrary
$\tilde a:L^{\otimes2}\dotimes{L,\iota_1}M\to
L^{\otimes2}\dotimes{L,\iota_2}M$. By the universal property of the tensor
product \cite[Proposition 7.3.4]{Li1}, specifying a pair of homomorphisms of
commutative $K$-algebras $f,g:L\to E$ is equivalent to specifying a
homomorphism $f\otimes g:L\dotimes{K}L\to E$. Thus $\tilde a$ defines the
following $E$-module homomorphism:
% segment-id: o014.li2.chapter7.module-descent.q010
\begin{equation*}
	\tilde{a}_{gf} := E \dotimes{f \otimes g} \tilde{a}: E \dotimes{L, f} M \to E \dotimes{L, g} M.
\end{equation*}

% segment-id: o014.li2.chapter7.module-descent.p010
As $(E,f,g)$ varies, we want to analyze the compatibility among these
homomorphisms. The formula above shows that
$(E,f,g):=(L\dotimes{K}L,\iota_1,\iota_2)$ is ``universal'' for this problem
and corresponds to $\tilde a$, while a general $(E,f,g)$ is obtained by
applying $E\dotimes{f\otimes g}(\mathord\cdot)$ to it. To analyze the
``universal'' case, for $1\leq i\neq j\leq3$ write:
% segment-id: o014.li2.chapter7.module-descent.q011
\begin{center}\begin{tabular}{|c|c|} \hline
		$\iota^3_{ij}: L^{\otimes 2} \to L^{\otimes 3}$ & insertion into tensor slots $(i,j)$ \\
		$\iota^3_i: L \to L^{\otimes 3}$ & insertion into the $i$th tensor slot \\ \hline
\end{tabular}\end{center}

% segment-id: o014.li2.chapter7.module-descent.p011
Using \eqref{eqn:flat-descent-map-aux}, define the $L^{\otimes3}$-module
homomorphism
% segment-id: o014.li2.chapter7.module-descent.q012
\[ \tilde{a}_{ji} := L^{\otimes 3} \dotimes{L, \iota^3_{ij}}\tilde{a} : \; L^{\otimes 3} \dotimes{L, \iota^3_i} M \to L^{\otimes 3} \dotimes{L, \iota^3_j} M. \]
Here we use $\iota^3_{ij}\iota_1=\iota^3_i$ and
$\iota^3_{ij}\iota_2=\iota^3_j$. When $M$ comes from a $K$-module, we may
identify
% segment-id: o014.li2.chapter7.module-descent.q013
\[ \tilde{a} \;\text{with}\; \identity_{L^{\otimes 2} \dotimes{K} N}, \quad \tilde{a}_{ji} \;\text{with}\; \identity_{L^{\otimes 3} \dotimes{K} N}. \]
Then $\tilde a_{31}=\tilde a_{32}\tilde a_{21}$ is clear. This is the desired
compatibility condition. Its important consequence is
% segment-id: o014.li2.chapter7.module-descent.q014
\begin{equation}\label{eqn:cocycle-condition-descent}
	\tilde{a}_{31} = \tilde{a}_{32} \tilde{a}_{21} \implies
	\begin{array}{l}
		\forall f, g, h: L \xrightarrow{K\text{-algebra}} E, \\
		\tilde{a}_{hf} = \tilde{a}_{hg} \tilde{a}_{gf};
	\end{array}
\end{equation}
the left-hand side is the ``universal version'' of the right-hand side.

% segment-id: o014.li2.chapter7.module-descent.d002
\begin{definition}[Descent datum]
	\index{descent datum}
	Define $\mathrm{Desc}_{K\to L}$ to be the following category. Its objects
	are data $(M,\tilde a)$, where $M$ is an $L$-module and
	$\tilde a:L^{\otimes2}\dotimes{L,\iota_1}M\to
	L^{\otimes2}\dotimes{L,\iota_2}M$ satisfies
	\begin{description}
		\item[Isomorphism condition] $\tilde a$ is an isomorphism of
		$L^{\otimes2}$-modules;
		\item[Cocycle condition] $\tilde a_{31}=\tilde a_{32}\tilde a_{21}$.
	\end{description}

	A morphism from $(M,\tilde a)$ to $(M',\tilde a')$ is an $L$-module
	homomorphism $f:M\to M'$ satisfying
	$(\identity^{\otimes2}\otimes f)\tilde a=
	\tilde a'(\identity^{\otimes2}\otimes f)$. Such data $(M,\tilde a)$ are
	called \emph{descent data}.
\end{definition}

% segment-id: o014.li2.chapter7.module-descent.p012
The discussion before \eqref{eqn:cocycle-condition-descent} shows that
$L\dotimes{K}(\mathord\cdot)$ gives a functor
$K\dcate{Mod}\to\mathrm{Desc}_{K\to L}$. We now relate descent data to
$L\dcate{Mod}^{\mathbf{L}}$.

% segment-id: o014.li2.chapter7.module-descent.le001
\begin{lemma}\label{prop:descent-data-vs-comodule}
	There is an isomorphism of categories from $L\dcate{Mod}^{\mathbf{L}}$ to
	$\mathrm{Desc}_{K\to L}$, given on objects by
	$(M,a)\leftrightarrow(M,\tilde a)$ and defined on morphisms in the
	standard way.

	Moreover, the isomorphism condition in the definition of
	$\mathrm{Desc}_{K\to L}$ may be replaced by the following condition:
	define $m:L^{\otimes2}\to L$ by
	$m(\ell\otimes\ell')=\ell\ell'$. The endomorphism $M\to M$ given by
	$L\dotimes{L^{\otimes2},m}\tilde a$ is $\identity_M$.
\end{lemma}
% segment-id: o014.li2.chapter7.module-descent.pf002
\begin{proof}
	Specifying $\tilde a$ is known to be equivalent to specifying
	$a:M\to L\dotimes{K}M$. The issue is to match the conditions on the
	descent datum $\tilde a$ with the conditions on $a$ that make $(M,a)$ a
	comodule. This is where the second assertion will be used.

	The comodule conditions are the coassociativity and counit laws. For the
	first, the equivalence of $\tilde a_{31}=\tilde a_{32}\tilde a_{21}$ with
	coassociativity is a mechanical verification left to the reader. For the
	second, write $\varphi$ for the composite
	$M\xrightarrow{a}L\dotimes{K}M\xrightarrow{\epsilon_M}M$. The counit law
	is equivalent to $\varphi=\identity$. Recalling that
	$\epsilon_M(\ell\otimes m)=\ell m$, a similar direct verification gives
	% segment-id: o014.li2.chapter7.module-descent.q015
	\[ \varphi = L \dotimes{L^{\otimes 2}, m} \tilde{a}. \]
	We next show that, under the assumption
	$\tilde a_{31}=\tilde a_{32}\tilde a_{21}$, the map $\tilde a$ is an
	isomorphism if and only if $\varphi=\identity_M$. This will complete the
	proof.

	Suppose $\tilde a$ is an isomorphism. Then $\varphi$ is also an
	isomorphism. In \eqref{eqn:cocycle-condition-descent}, take $E=L$ and
	$f=g=h=\identity_L$. The right-hand side readily becomes
	$\varphi=\varphi^2$, and hence $\varphi=\identity_M$.

	Conversely, suppose $\varphi=\identity_M$. For every homomorphism of
	$K$-algebras $f:L\to E$, change scalars on $\tilde a$ along the two paths
	of the commutative diagram
	% segment-id: o014.li2.chapter7.module-descent.q016
	\[\begin{tikzcd}
		L^{\otimes 2} \arrow[rd, "{f \otimes f}"'] \arrow[r, "m"] & L \arrow[d, "f"] \\
		& E.
	\end{tikzcd}\]
	This shows that $\tilde a_{ff}=\identity$. Together with
	\eqref{eqn:cocycle-condition-descent}, it follows that
	$\tilde a_{fg}^{-1}=\tilde a_{gf}$ for every $(E,f,g)$. In particular,
	the ``universal version'' $\tilde a$ is also an isomorphism. This proves
	the lemma.
\end{proof}

% segment-id: o014.li2.chapter7.module-descent.th002
\begin{theorem}[Second form of flat descent]\label{prop:flat-descent-2}
	Let $L$ be a faithfully flat commutative $K$-algebra. Then
	% segment-id: o014.li2.chapter7.module-descent.q017
	\[ L \dotimes{K} (\cdot): K\dcate{Mod} \to \mathrm{Desc}_{K \to L} \]
	is an equivalence of categories. One of its quasi-inverses,
	$\tilde{\mathbb{L}}$, may be taken to be the equalizer
	% segment-id: o014.li2.chapter7.module-descent.q018
	\[\begin{tikzcd}
		0 \arrow[r] & \tilde{\mathbb{L}}(M, \tilde{a}) \arrow[r] & M \arrow[shift left, r, "a"] \arrow[shift right, r, "{m \mapsto 1 \otimes m}"'] & L \dotimes{K} M,
	\end{tikzcd}\]
	where $a$ is the map corresponding to $\tilde a$ under
	\eqref{eqn:flat-descent-map-aux}.
\end{theorem}
% segment-id: o014.li2.chapter7.module-descent.pf003
\begin{proof}
	By Lemma~\ref{prop:descent-data-vs-comodule}, this is merely a
	reformulation of the flat descent Theorem~\ref{prop:flat-descent}.
\end{proof}

% segment-id: o014.li2.chapter7.module-descent.p013
Theorem~\ref{prop:flat-descent-2} is the usual formulation of flat descent in
algebraic geometry. There a faithfully flat homomorphism $K\to L$
corresponds, in a certain sense, to a covering of a geometric object.

% segment-id: o014.li2.chapter7.module-descent.p014
The datum $\tilde a$, or the corresponding $a$, is sometimes also expressed
as a homomorphism of left $L\dotimes{K}L$-modules, that is, of
$(L,L)$-bimodules,
% segment-id: o014.li2.chapter7.module-descent.q019
\begin{align*}
	\overline{a}: M \dotimes{K} L & \to L \dotimes{K} M \\
	m \otimes \ell & \mapsto \sum_{i=1}^n \ell_i \otimes \ell m_i,
\end{align*}
still assuming $a(m)=\sum_{i=1}^n\ell_i\otimes m_i$. This is obtained by
contracting the tensor products in \eqref{eqn:flat-descent-map-aux} in the
appropriate way and has the advantage of a more compact form. A contracted
version for $\tilde a_{ji}$ is obtained similarly.

% segment-id: o014.li2.chapter7.module-descent.r002
\begin{remark}[Flat descent for algebras and modules]\label{rem:flat-descent-alg}
	We can go on to discuss how $L$-algebras, or modules over those algebras,
	descend to $K$. Concretely, suppose
	$(M,a),(M',a')\in\Obj(L\dcate{Mod}^{\mathbf{L}})$ are the images of
	$N,N'\in\Obj(K\dcate{Mod})$. Recall that
	$L\dotimes{K}(\mathord\cdot):K\dcate{Mod}\to L\dcate{Mod}$ is a
	monoidal functor \cite[Proposition 6.6.10]{Li1}; that is, there are
	canonical isomorphisms such as
	% segment-id: o014.li2.chapter7.module-descent.q020
	\begin{equation}\label{eqn:change-of-ring-monoidal}\begin{aligned}
			(L \dotimes{K} N_1) \dotimes{L} (L \dotimes{K} N_2) & \rightiso L \dotimes{K} (N_1 \dotimes{K} N_2) \\
			(\ell_1 \otimes x) \otimes (\ell_2 \otimes y) & \mapsto \ell_1 \ell_2 \otimes (x \otimes y).
	\end{aligned}\end{equation}
	Use this to give $M\dotimes{L}M'$ the $\mathbf{L}$-comodule structure
	coming from $N\dotimes{K}N'$. Now take $N'=N$. Giving $N$ a
	$K$-algebra structure is equivalent to giving $M$ an $L$-algebra
	structure such that the multiplication $M\dotimes{L}M\to M$ and unit
	$L\to M$ are both morphisms in $L\dcate{Mod}^{\mathbf{L}}$; all
	properties such as associativity required of a $K$-algebra come from $M$.
	The descent of $M$-modules to $N$-modules should be understood in the same
	way. For this purpose, descent data of the form
	\eqref{eqn:flat-descent-map-aux} may be more convenient than comodules over
	$\mathbf{L}$, because all the functors involved are monoidal on $M$,
	whereas $\mathbf{L}$ itself is not a monoidal functor.
\end{remark}
